\documentclass[a4paper,12pt]{article}
\usepackage[utf8]{inputenc}
\usepackage[russian]{babel}
\usepackage{graphicx}
\usepackage{hyperref}
\usepackage{minted}
\usemintedstyle{friendly}
\usepackage{xcolor}
\usepackage{geometry}
\usepackage{amsmath}
\usepackage{enumitem}
\usepackage{float}
\usepackage{indentfirst}
\usepackage{caption}
\usepackage{amsfonts}
\usepackage{tikz}
\usetikzlibrary{shapes.geometric, arrows.meta, positioning}

\documentclass[14pt]{article}
\usepackage[utf8]{inputenc}
\usepackage[russian]{babel}
\usepackage{geometry}
\usepackage{minted}

\geometry{top=2cm, bottom=2cm, left=2.5cm, right=1.5cm}

% Настройки для листингов кода
\setminted{
    fontsize=\small,
    breaklines=true,
    autogobble=true,
    frame=single,
}

\begin{document}

\begin{titlepage}
    \centering
    % Верхний блок
    {\bfseries МИНИСТЕРСТВО НАУКИ И ВЫСШЕГО ОБРАЗОВАНИЯ РФ \\
    ФГБОУ ВО «Белгородский государственный технологический университет им. В.Г. Шухова»\par}
    
    \vfill
    
    % Центральный блок
    {\bfseries\Large Практическое задание №9\par}
    \vspace{0.5cm}
    {\large по дисциплине: Основы программирования и алгоритмизации \par}
    {\large тема: «Линейная алегбра и сжатие информации» \par}
    
    \vfill
    
    % Блок авторов
    \begin{flushright}
    \begin{tabular}{l}
        \large Выполнил: студент группы ВТ-252 \\
        \large Дронь Дмитрий Алексеевич \\
        \large Проверила: \\
        \large Стеценко Оксана Николаевна
    \end{tabular}
    \end{flushright}
    
    \vfill
    
    % Нижний блок
    {\large Белгород, 2026 г. \par}
\end{titlepage}

\newpage
\begin{center}
\section*{Цель работы}
\end{center}
Цель методических указаний организовать самостоятельную работу студента по выполнению индивидуальных заданий, направленных на закрепление навыков программирования на языке высокого уровня и способов решения математических задач с применением программы для работы с электронными таблицами MS Excel.
\newpage
\begin{center}
\section*{Индивидуальное задание №9}
\end{center}
\textbf{Задача 9.} Матрицу A(m,n) разложить на слагаемые: A = B + C + D, где B — строго нижнетреугольная, C — диагональная, D — строго верхнетреугольная матрицы того же размера. Для экономии памяти матрицы В, С, D хранить в упакованном виде в одномерных массивах; вывести все треугольные матрицы в «треугольном» виде.
\newpage
\begin{center}
\section*{Блок-схема}
\scalebox{0.60}{
\begin{tikzpicture}[
    >=Stealth,
    line/.style={draw,-Stealth,thick},
    block/.style={
        rectangle,
        rounded corners,
        draw,
        align=center,
        text width=3.5cm,
        minimum height=1cm
    },
    blockplain/.style={
        rectangle,
        draw,
        align=center,
        text width=3.5cm,
        minimum height=1cm
    },
    decision/.style={
        diamond,
        draw,
        align=center,
        aspect=2.2,
        text width=2.4cm,
        inner sep=2pt
    }
]

%------------------ Узлы ------------------
\node[block] (input) at (0,0)
{
decompose\\
n, A[MAXN][MAXN]\\
B[], C[], D[]
};

\node[blockplain] (init) at (0,-2.8)
{
idxB = 0\\
idxC = 0\\
idxD = 0\\
i = 0
};

\node[decision] (icheck) at (0,-6)
{i < n};

\node[blockplain] (jzero) at (0,-8.5)
{j = 0};

\node[decision] (jcheck) at (0,-11)
{j < n};

\node[decision] (igtj) at (0,-14)
{i > j};

\node[blockplain] (storeB) at (0,-17)
{
B[idxB] = A[i][j];\\
idxB++
};

\node[decision] (ieqj) at (0,-20)
{i == j};

\node[blockplain] (storeC) at (0,-23)
{
C[idxC] = A[i][j];\\
idxC++
};

\node[decision] (iltj) at (0,-26)
{i < j};

\node[blockplain] (storeD) at (0,-29)
{
D[idxD] = A[i][j];\\
idxD++
};

\node[blockplain] (jinc) at (0,-31.8)
{j++};

\node[blockplain] (iinc) at (0,-34.5)
{i++};

\node[block] (ret) at (0,-37.2)
{возврат r};

%------------------ Стрелки ------------------
\draw[] (input) -- (init);
\draw[] (init) -- (icheck);

% icheck: jzero
\draw[]
    (icheck.south) -- ++(0,-0.2)
    -- (jzero.north);

% icheck: ret (влево и вниз)
\draw[line]
    (icheck.west) -- ++(-3.8,0)
    -- ++(0,-30)
    -- ++(5.55,0)
    -- (ret.north);

\draw[] (jzero.south) -- (jcheck.north);

% jcheck: igtj
\draw[]
    (jcheck.south) -- ++(0,-0.2)
    -- (igtj.north);

% jcheck: нет -> iinc
\draw[line]
    (jcheck.west) -- ++(-2,0)
    -- ++(0,-22)
    -- ++(3.825,0)
    -- (iinc.north);

% igtj: storeB
\draw[]
    (igtj.south) -- ++(0,-0.2)
    -- (storeB.north);


\draw[] (storeB.south) -- (ieqj.north);

% ieqj: storeC
\draw[]
    (ieqj.south) -- ++(0,-0.2)
    -- (storeC.north);


\draw[] (storeC.south) -- (iltj.north);

% iltj: storeD
\draw[]
    (iltj.south) -- ++(0,-0.2)
    -- (storeD.north);

\draw[] (storeD.south) -- (jinc.north);

% jinc -> jcheck
\draw[line]
    (jinc.east) -- ++(2.5,0)
    -- ++(0,22.3)
    -- ++(-4.4,0)
    -- (jcheck.north);

% iinc -> icheck петля
\draw[line]
    (iinc.east) -- ++(4,0)
    -- ++(0,30)
    -- ++(-5.9,0)
    -- (icheck.north);

\end{tikzpicture}
}
\end{center}
\newpage
\begin{center}
\scalebox{0.60}{
\begin{tikzpicture}[
    >=Stealth,
    line/.style={draw,-Stealth,thick},
    block/.style={
        rectangle,
        rounded corners,
        draw,
        align=center,
        text width=3.6cm,
        minimum height=1.2cm
    },
    blockplain/.style={
        rectangle,
        draw,
        align=center,
        text width=2.6cm,
        minimum height=1cm
    },
    io/.style={
        trapezium,
        trapezium left angle=70,
        trapezium right angle=110,
        trapezium stretches=true,
        draw,
        align=center,
        text width=2.2cm,
        minimum height=1cm
    },
    decision/.style={
        diamond,
        draw,
        align=center,
        aspect=2.0,
        text width=2.4cm,
        inner sep=1pt
    }
]

%------------------ Узлы ------------------
\node[block] (input) at (0,0)
{
output\_triangular\\
n, arr[], type
};

\node[blockplain] (init) at (0,-2.2)
{i = 0};

\node[decision] (icheck) at (0,-4.5)
{i < n};

\node[blockplain] (jzero) at (0,-6.8)
{j = 0};

\node[decision] (jcheck) at (0,-9.1)
{j < n};

% --- Блок 1: type == 1 && i > j ---
\node[decision] (cond1) at (0,-11.8)
{
type == 1\\
\&\& i > j
};

\node[io] (out1_no) at (-3.2,-13.6)
{0};

\node[io] (out1_yes) at (3.2,-13.6)
{arr[idx]};

\node[blockplain] (inc1) at (3.2,-15.6)
{idx++};

\coordinate (merge1) at (0,-17.0);

% --- Блок 2: type == 2 && i == j ---
\node[decision] (cond2) at (0,-19.0)
{
type\\
== 2 \&\& i ==\\
j
};

\node[io] (out2_no) at (-4.2,-20.8)
{0};

\node[io] (out2_yes) at (4.2,-20.8)
{arr[idx]};

\node[blockplain] (inc2) at (4.2,-22.8)
{idx++};

\coordinate (merge2) at (0,-24.2);

% --- Блок 3: type == 3 && i < j ---
\node[decision] (cond3) at (0,-26.2)
{
type == 3\\
\&\& i < j
};

\node[io] (out3_no) at (-3.2,-28.0)
{0};

\node[io] (out3_yes) at (3.2,-28.0)
{arr[idx]};

\node[blockplain] (inc3) at (3.2,-30.0)
{idx++};

\coordinate (merge3) at (0,-31.4);

% --- Инкременты циклов ---
\node[blockplain] (jinc) at (0,-33.0)
{j++};

\node[blockplain] (iinc) at (0,-36.0)
{i++};

%------------------ Стрелки и линии ------------------
\draw[] (input.south) -- (init.north);
\draw[line] (init.south) -- (icheck.north);
\draw[] (icheck.south) -- (jzero.north);
\draw[line] (jzero.south) -- (jcheck.north);
\draw[] (jcheck.south) -- (cond1.north);

% Ветвление 1
\draw[] (cond1.west) -| (out1_no.north) node[pos=0.25, above, font=\small] {нет};
\draw[] (cond1.east) -| (out1_yes.north) node[pos=0.25, above, font=\small] {да};
\draw[] (out1_yes.south) -- (inc1.north);
\draw[] (out1_no.south) |- (merge1);
\draw[] (inc1.south) |- (merge1);
\draw[] (merge1) -- (cond2.north);

% Ветвление 2
\draw[] (cond2.west) -| (out2_no.north) node[pos=0.25, above, font=\small] {нет};
\draw[] (cond2.east) -| (out2_yes.north) node[pos=0.25, above, font=\small] {да};
\draw[] (out2_yes.south) -- (inc2.north);
\draw[] (out2_no.south) |- (merge2);
\draw[] (inc2.south) |- (merge2);
\draw[] (merge2) -- (cond3.north);

% Ветвление 3
\draw[] (cond3.west) -| (out3_no.north) node[pos=0.25, above, font=\small] {нет};
\draw[] (cond3.east) -| (out3_yes.north) node[pos=0.25, above, font=\small] {да};
\draw[] (out3_yes.south) -- (inc3.north);
\draw[] (out3_no.south) |- (merge3);
\draw[] (inc3.south) |- (merge3);
\draw[] (merge3) -- (jinc.north);

% Выход из вложенного цикла (j < n -> i++)
\draw[line]
    (jcheck.west) -- ++(-5.0,0)
    |- ([yshift=0.6cm]iinc.north)
    -- (iinc.north);

% Петля j++ -> j < n
\draw[line]
    (jinc.east) -- ++(5.0,0)
    |- ([yshift=0.5cm]jcheck.north)
    -- (jcheck.north);

% Петля i++ -> i < n
\draw[line]
    (iinc.east) -- ++(7.0,0)
    |- ([yshift=0.5cm]icheck.north)
    -- (icheck.north);

\end{tikzpicture}
}
\end{center}

\newpage
\begin{center}
\section*{Программа}
\end{center}
\begin{minted}{c}
/* Летняя практика. Индивидуальное задание №9 «Линейная алгебра и сжатие информации»
   Выполнил студент группы ВТ-252 Дронь Д.А.
   Задание 9: Матрицу A(m,n) разложить на слагаемые: A = B + C + D, где B — строго нижнетреугольная, C — диагональная, D — строго верхнетреугольная матрицы того же размера. Для экономии памяти матрицы В, С, D хранить в упакованном виде в одномерных массивах; вывести все треугольные матрицы в «треугольном» виде.
*/
#include <stdio.h>

#define MAXN 50

// Структура результата
typedef struct {
    int error_flag;
} Result;

// Ввод матрицы A(n,n)
int input_matrix(int *n, int A[MAXN][MAXN]) {
    int flag = 0;
    int tn;

    printf("Введите размер матрицы n: ");
    if (scanf("%d", &tn) != 1) {
        flag = 1;
    }

    if (flag == 0) {
        if (tn <= 0 || tn > MAXN) {
            flag = 2;
        }
    }

    if (flag == 0) {
        *n = tn;
        printf("Введите матрицу %d x %d:\n", tn, tn);

        int i = 0;
        while (i < tn && flag == 0) {
            int j = 0;
            while (j < tn && flag == 0) {
                if (scanf("%d", &A[i][j]) != 1) {
                    flag = 1;
                }
                j++;
            }
            i++;
        }
    }

    return flag;
}

/*
Разложение A = B + C + D
B — строго нижнетреугольная (i > j)
C — диагональная (i == j)
D — строго верхнетреугольная (i < j)
Все три матрицы храним в одномерных массивах
*/
Result decompose(int n, int A[MAXN][MAXN],
                 int B[], int C[], int D[])
{
    Result r;
    r.error_flag = 0;

    int idxB = 0;
    int idxC = 0;
    int idxD = 0;

    int i = 0;
    while (i < n) {
        int j = 0;
        while (j < n) {

            // Строго нижний треугольник
            if (i > j) {
                // сохраняем A[i][j] в B
                B[idxB] = A[i][j];
                idxB++;
            }

            // Диагональ
            if (i == j) {
                // сохраняем A[i][j] в C
                C[idxC] = A[i][j];
                idxC++;
            }

            // Строго верхний треугольник
            if (i < j) {
                // сохраняем A[i][j] в D
                D[idxD] = A[i][j];
                idxD++;
            }

            j++;
        }
        i++;
    }

    return r;
}

/*
Вывод треугольной матрицы из упакованного массива
Тип:
 type = 1 → нижняя
 type = 2 → диагональ
 type = 3 → верхняя
*/
void output_triangular(int n, int arr[], int type) {
    int idx = 0;

    printf("\n");

    int i = 0;
    while (i < n) {
        int j = 0;
        while (j < n) {

            if (type == 1) { // нижняя
                if (i > j) {
                    printf("%d ", arr[idx]);
                    idx++;
                } else {
                    printf("0 ");
                }
            }

            if (type == 2) { // диагональ
                if (i == j) {
                    printf("%d ", arr[idx]);
                    idx++;
                } else {
                    printf("0 ");
                }
            }

            if (type == 3) { // верхняя
                if (i < j) {
                    printf("%d ", arr[idx]);
                    idx++;
                } else {
                    printf("0 ");
                }
            }

            j++;
        }
        printf("\n");
        i++;
    }
}

int main(void) {
    int n;
    int A[MAXN][MAXN];

    int err = input_matrix(&n, A);

    switch (err) {
        case 0:
            break;
        case 1:
            printf("Ошибка: некорректный формат ввода.\n");
            return 1;
        case 2:
            printf("Ошибка: n должно быть в диапазоне 1..%d.\n", MAXN);
            return 1;
        default:
            printf("Неизвестная ошибка ввода.\n");
            return 1;
    }

    // Размеры упакованных массивов
    int sizeB = n * (n - 1) / 2;
    int sizeC = n;
    int sizeD = n * (n - 1) / 2;

    int B[sizeB];
    int C[sizeC];
    int D[sizeD];

    Result r = decompose(n, A, B, C, D);

    switch (r.error_flag) {
        case 0:
            printf("\nСтрого нижняя треугольная матрица B:\n");
            output_triangular(n, B, 1);

            printf("\nДиагональная матрица C:\n");
            output_triangular(n, C, 2);

            printf("\nСтрого верхняя треугольная матрица D:\n");
            output_triangular(n, D, 3);
            break;

        default:
            printf("Неизвестная ошибка вычислений.\n");
            break;
    }

    return 0;
}


\end{minted}
\newpage
\begin{center}
\section*{Тесты}
\end{center}
\begin{figure}[H]
    \centering
    \includegraphics[width=0.9\textwidth]{test.png}
    \caption{Результат выполнения программы}
\end{figure}
\newpage
\begin{center}
\section*{Вывод}
\end{center}
В ходе выполнения работы мне удалось организовать самостоятельную работу по выполнению индивидуальных заданий, направленных на закрепление навыков программирования на языке высокого уровня и способов решения математических задач с применением программы для работы с электронными таблицами MS Excel.

\end{document}




